\begin{figure}[!htbp]
\centering
\pgfplotstableread{
x y em ep fam res
1.1096 0 0.0608 0.0688 1 1
1.0004 1 0.0006 0.0006 1 1
0.9979 2 0.0013 0.0013 1 1
0.4458 3 0.0143 0.0145 1 1
0.6159 4 0.0247 0.0266 1 1
0.9799 5 0.0438 0.0408 1 0
0.9894 6 0.0133 0.0125 0 0
0.9667 7 0.0266 0.0300 0 0
0.9509 8 0.0334 0.0372 0 0
0.9528 9 0.0290 0.0280 0 1
0.8358 10 0.0502 0.0462 0 1
}\forestdata
\pgfplotsset{forest ticks/.style={ytick={0,1,2,3,4,5,6,7,8,9,10}, yticklabels={{H5: TimesFM-3 uni vs Naive, FX and FRED},{H4d: multivariate vs univariate, spurious covariate},{H4c: multivariate vs univariate, independent series},{H4b: multivariate vs univariate, known covariate},{H4a: multivariate vs univariate, lead-lag 24},{H2: W-seen vs W-post DiD, TimesFM-3 vs Chronos-2},{P5: TimesFM-3 eval vs TimesFM-3},{P4a: TimesFM-3 vs TimesFM-3 uni},{P3: TimesFM-3 uni vs TimesFM-2.5},{P2: TimesFM-3 uni vs Chronos-2 uni},{P1: TimesFM-3 uni vs Stat.\ ensemble}}, ymin=-0.7, ymax=10.7}}

\begin{tikzpicture}
\begin{axis}[
  width=0.62\textwidth, height=9.2cm,
  xmode=log, log basis x=10, xmin=0.4, xmax=1.3,
  xtick={0.5,0.6,0.7,0.8,0.9,1,1.1,1.2}, xticklabels={0.5,0.6,0.7,0.8,0.9,1,1.1,1.2},
  forest ticks, yticklabel style={font=\scriptsize, align=right},
  xlabel={ratio of errors, A / B (log scale; $<1$: A better)}, xlabel style={font=\footnotesize},
  tick label style={font=\footnotesize},
  axis lines*=left, ymajorgrids=false, xmajorgrids=true, grid style={black!10},
  clip=false]
\addplot[black!40, dashed] coordinates {(1,-0.7) (1,10.7)};
\addplot[only marks, mark=*, mark size=1.6pt, black,
  error bars/.cd, x dir=both, x explicit, error bar style={line width=0.8pt}]
  table[x=x, y=y, x error minus=em, x error plus=ep] {\forestdata};
\end{axis}
\end{tikzpicture}
\caption{Primary (family P, MASE) and hypothesis (family H) contrasts: estimate and two-way cluster bootstrap 95\% interval
(unadjusted). P-tests are on post-release Tier B data: P1--P3 without covariates on the same series and windows (context lengths differ on
hourly data), P4a on the group tasks (no covariates), P5 between two configurations with the covariates of load and river flow; H4a--H4d are on controlled
multivariate data; H5 is on exchange rates, Treasury yields and oil prices. H3 is omitted (its cell never activates the mechanism
it tests, Section~\ref{sec:controlled}). Decisions use Holm-adjusted $p$-values (Tables~\ref{tab:familyp} and~\ref{tab:familyh}).}
\label{fig:forest}
\end{figure}
